\section{Support and extremal corner obstructions}\label{sec:support}
Let $M(k)$ be the least number of cubes guaranteeing a corner solution when
the collection contains at most $k$ varieties. We determine these thresholds
and classify all fourteen-cube obstructions into 56 symmetry classes. Write $K(X)$ for the set of the eight distinct cyclic corner types of $X$.

\paragraph{Directed corners.}
Write $V=\col$ and encode a tricoloured cyclic triple $(i,j,k)$ by the complementary
ordered pair $ab$ for which $(a,b,i,j,k)$ is an even permutation of $V$.
Cyclically rotating the triple preserves $ab$; reversing it gives $ba$.
For example, $013\mapsto42$ and $031\mapsto24$.
Write $b(r,p)$ for the bicoloured corner $rrp$.
The tricoloured corners of an adjacent-repeat variety form a cut
\[
 C(S)=S\times(V\setminus S),\qquad |S|\in\{2,3\}.
\]
Its smaller side is a pair $P$, its repeated colour is $r\notin P$, and
its other corners are $b(r,p)$ for $p\in P$. The pair, repeated colour and
sign determine a unique adjacent-repeat variety. Call $C(P)$ positive and
$C(V\setminus P)$ negative. For example, $(0,1,0,2,3,4)$ has cut $C(34)$
and bicoloured corners $b(0,3),b(0,4)$. Colour permutations give all adjacent
varieties; odd permutations also reverse the arcs. A string such as $34$
in $C(34)$ denotes $\{3,4\}$.
An opposite-repeat variety with singleton opposite pairs $P,Q$ has arcs
$(P\times Q)\cup(Q\times P)$, a bidirected $K_{2,2}$ omitting its repeated
colour.

The identity
\[
 C(S)\cap C(T)=(S\cap T)\times\bigl(V\setminus(S\cup T)\bigr)
\]
gives the following tricoloured intersection sizes for smaller pairs $P,Q$:
\[
\begin{array}{c|ccc}
 |P\cap Q|&2&1&0\\\hline
 \text{same sign}&6&2&0\\
 \text{opposite signs}&0&1&4
\end{array}
\]
There are additionally $|P\cap Q|$ shared bicoloured corners when the
repeated colours agree, and none otherwise. In particular, a singleton
full intersection is tricoloured and joins opposite signs with different
repeated colours. By Lemma~\ref{lem:cube-incidence}, a pair involving an
opposite-repeat variety shares at least two corners; two opposite varieties
sharing exactly two corners share one bidirected edge.

\begin{lemma}\label{lem:disjoint-pair}
Four copies of $A$, together with any four sources individually
corner-disjoint from $A$, can build an opposite-repeat target.
\end{lemma}
\begin{proof}
The anchor and all four other sources are adjacent-repeat. Reflect the entire instance
if necessary so that $A$ has cut $C(P)$, $|P|=2$, and put $R=V\setminus P$.
Each disjoint source has cut $C(R)$ or $C(Q)$ for a pair $Q\subset R$.
Choose $Q$ maximizing frequency among all three pair-cut classes,
including those of frequency zero. Four $A$ copies fill
$P\times Q$ of the opposite target with singleton pairs $P,Q$.
On $Q\times P$, sources of classes $C(Q),C(R)$ are unrestricted; each other
pair class fits one two-position row. Neither row-only class occurs more
than twice, by maximality of the chosen class. Place them in their rows
and fill the remaining positions with the unrestricted sources.
\end{proof}

Consequently, eleven sources containing seven copies of one variety solve:
one of the other four sources sharing a corner completes the seven-copy
target; otherwise all four are disjoint and the lemma applies.

\begin{samepage}
\begin{theorem}[Support thresholds]\label{thm:support}
For every integer $k\ge1$,
\[
 M(k)=\begin{cases}8&k=1,\\13&k=2,\\15&k\ge3.\end{cases}
\]
Three is therefore the minimum number of varieties in a fourteen-cube obstruction.
\end{theorem}
\end{samepage}
\begin{proof}
The one-variety case is immediate. Thirteen sources in two varieties include
seven of one variety, so the preceding consequence gives the upper bound.
For the lower bound take six copies each of the original $B,C$ from
Lemma~\ref{lem:lower}. Their corner union supports only $B,C$: removing
$A$ from that lemma's three-variety union cannot introduce another target,
and $A$ has five corners absent from both. The sole shared corner is $042$,
so either remaining target requires seven copies of its own variety.
Thus $6B+6C$ fails. For $k\ge3$, Theorem~\ref{thm:main} and
Lemma~\ref{lem:lower} give the matching bounds.
\end{proof}

\begin{theorem}[Multiplicity bound]\label{thm:multiplicity-five}
Every fourteen-cube obstruction has maximum multiplicity at most five.
\end{theorem}
\begin{proof}
The seven-copy case follows above. Suppose there are exactly six copies
of $A$ and eight other sources. Two others fitting distinct $A$ corners
complete $A$; four disjoint others invoke Lemma~\ref{lem:disjoint-pair}.
In a failure, therefore, at most three others are disjoint and at least
five meet $A$, all in the same singleton corner $k$.

The singleton intersection implies that $A$ is adjacent-repeat, since
an opposite-repeat source shares at least two corners with every variety.
Normalize the cut of $A$ to $C(P)$ with $P=\{p,p'\}$, and put $R=V\setminus P$.
The intersection table makes $k=pq$ for some $q\in R$ and forces every
meeting source to have cut
\[
 C(T),\qquad T=\{p\}\cup(R\setminus\{q\}).
\]
There are at most three varieties with this common cut. Reserve two
identical $B$ copies for its bicoloured corners, set aside any three further
common-cut copies, and reserve one $A$ copy for $k$. A disjoint source's
intersection with $C(T)$ is
\[
\begin{array}{c|c}
 \text{source cut}&\text{intersection with }C(T)\\\hline
 C(R)&(R\setminus\{q\})\times\{p'\}\\
 C(R\setminus\{q\})&(R\setminus\{q\})\times\{p',q\}\\
 C(\{q,r\}),\ r\in R\setminus\{q\}&\{rp'\}
\end{array}
\]
and always avoids $k$. Exactly three of the eight non-$A$ copies remain.
Each meets $C(T)\setminus\{k\}$: a common-cut copy fits all five positions,
and the displayed intersections for disjoint sources are nonempty.
If two remaining copies fit distinct corners there, use them and fill the
other three positions with the set-aside common-cut copies. Together with
the two reserved $B$ copies and the $A$ at $k$, this completes $B$.

Otherwise the three nonempty compatibility sets are one common singleton.
Indeed, a set containing two corners can be paired at distinct corners
with any other nonempty set; two different singleton sets also give such
a pair. None of the three remaining copies therefore has the common cut.
They are all disjoint from $A$, and the table identifies their shared singleton
as $rp'$. They hence have cut $C(Q)$ with $Q=\{q,r\}$. Retarget to the opposite
variety with singleton pairs $P,Q$: four $A$ copies fill $P\times Q$;
three disjoint sources and one $B$ at $rp'$ fill $Q\times P$.
Every case gives an assembly, independently of the fifteen-cube theorem.
\end{proof}

For classification, equivalence permits common colour relabelling and
simultaneous reflection of the entire collection. Reflection is not an
additional move for an individual source. In directed notation this is
arbitrary vertex relabelling and independent global arc reversal.

\paragraph{The generalized family.}
Let $A$ be the opposite variety with singleton pairs $\{1,2\},\{3,4\}$.
For $r\in\{2,3,4\}$, let $B_r$ have cut $C(01)$ and repeat $r$;
for $s\in\{1,2,4\}$, let $C_s$ have cut $C(124)$ and repeat $s$.
In the notation of Lemma~\ref{lem:lower}, $B_2=B$ and $C_1=C$.

\begin{samepage}
\begin{corollary}\label{cor:extremal-family}
With the same $A,B_r,C_s$, every collection
\[
 4A+\sum_{r\in\{2,3,4\}}x_rB_r+\sum_{s\in\{1,2,4\}}y_sC_s,
 \qquad \sum_r x_r=\sum_s y_s=5,
\]
where the counts are nonnegative integers, is an obstruction if and only if
$x_2y_2=x_4y_4=0$.
\end{corollary}
\end{samepage}
\begin{proof}
When $x_2y_2=x_4y_4=0$, the two groups have disjoint repeated-colour supports.
Their tricoloured union is contained in that of $A$, $C(01)$ and $C(124)$,
which misses the arcs $01,21,30,34$. Its bidirected edges are
$02,04,13,14,23,24$. No four-cycle passes through $0$, since its neighbours
$2,4$ have no second common neighbour; on the other four vertices the only
four-cycle is $A$. Thus $A$ is the only covered opposite-repeat target.
For an adjacent target, the bicoloured corners determine an active $B_r$ or
$C_s$, up to reversing its cut. The reversed cuts require absent arcs
$21,01$, respectively. Each remaining target has the following positions
private to its indicated source group:
\[
\begin{array}{c|l|c|c}
 \text{target}&\text{private positions}&\text{demand}&\text{supply}\\\hline
 A&24,31,32,41,42&5&4\\
 B_r&02,03,04,12\text{ and both bicoloured corners}&6&5\\
 C_s&10,20,40,43\text{ and both bicoloured corners}&6&5
\end{array}
\]
All targets therefore fail. Missing varieties within either group can only
remove suppliers.
Conversely, suppose both groups contain repeat $r\in\{2,4\}$. Build $B_r$:
reserve one $B_r$ at $b(r,1)$ and one $C_r$ at $b(r,0)$; the other four
$B$-group copies fill $02,03,04,12$, and two $A$ copies fill $13,14$.
These are eight distinct sources, giving a solution.
Taking $y_1=5$ gives a 21-member
subfamily. The general family also contains, for example,
$4A+2B_2+3B_3+2C_1+3C_4$, whose maximum multiplicity is four.
\end{proof}

% Staged section: include from extremal.tex only after the complete enumeration
% replay and physical-classification axiom audit have passed.
\subsection{The complete fourteen-cube classification}\label{sec:classification14}
The preceding family accounts for most, but not all, extremal obstructions.
There are two further forms, both using six varieties. Throughout this
subsection equivalence means common colour relabelling and optional
simultaneous reflection of the collection. A labelled collection keeps the
five colour names fixed; identical physical copies remain indistinguishable.

\paragraph{An all-adjacent family.}
Choose a colour $o$, an ordered partition $P,Q$ of $V\setminus\{o\}$ into
pairs, and a derangement $\rho$ of those four colours. Take one copy of each
repeat-$o$ variety with cuts $C(P)$ and $C(P\cup\{o\})$. For each $v\in P$,
take three copies with cut $C(\{v,o\})$ and repeat $\rho(v)$; for each
$v\in Q$, take three copies with cut $C(V\setminus\{v,o\})$ and repeat
$\rho(v)$. The condition $\rho(v)\ne v$ makes every variety well defined.
These collections have multiplicities $(3,3,3,3,1,1)$.

Their tricoloured union omits exactly $Q\times P$. Its bidirected graph is
the two triangles on $P\cup\{o\}$ and $Q\cup\{o\}$, sharing only $o$.
It has no four-cycle, so no opposite-repeat target is covered. For any
repeat other than $o$, the bicoloured corners fix the smaller pair of its
unique three-copy variety; the reverse cut requires a missing arc.
For repeat $o$, a smaller pair crossing $P,Q$ also requires a missing arc
in either orientation. Thus precisely the six source varieties are covered.

For $v\in P$, with $w$ the other member of $P$, the three-copy target's
private arcs are $vw,ow$. For $v\in Q$, with $w$ the other member of $Q$,
they are $wv,wo$. Its two bicoloured corners are also private, since $\rho$
is bijective. Each three-copy variety therefore faces four private positions, and each
singleton faces its two private bicoloured positions. Every target fails.

The colour $o$ is recovered from the two singleton varieties, and the
missing cut recovers $P,Q$ up to reversal. The eight-element partition
stabilizer acts on $\rho$ by conjugation. For $o=4$, $P=\{0,3\}$ and
$Q=\{1,2\}$, its four orbits are represented by
\[
\begin{array}{c|c|r}
 \text{type}&\rho&\text{labelled orbit size}\\\hline
 \text{within-pair transpositions}&(0\ 3)(1\ 2)&30\\
 \text{cross-pair transpositions}&(0\ 1)(3\ 2)&60\\
 \text{alternating four-cycle}&(0\ 1\ 3\ 2)&60\\
 \text{nonalternating four-cycle}&(0\ 3\ 1\ 2)&120
\end{array}
\]
The normalized orbit sizes are $1,2,2,4$, accounting for the nine
derangements. There are $5\binom42=30$ ordered templates, giving 270
labelled collections in these four classes.

\paragraph{A mixed exception.}
The remaining form is the following single collection, in the face order
$(+x,-x,+y,-y,+z,-z)$:
\[
\begin{array}{c|c|c|c}
 \text{variety}&\text{face tuple}&\text{copies}&\text{private positions}\\\hline
 X_0&(0,1,0,2,3,4)&5&6\\
 X_1&(0,4,1,2,3,1)&1&2\\
 X_2&(0,4,1,2,3,2)&1&2\\
 X_3&(0,4,1,3,3,2)&1&2\\
 O&(0,4,1,2,3,3)&3&4\\
 Y&(0,4,1,4,3,2)&3&4
\end{array}
\]
Only $O$ is opposite-repeat. The arc union misses $04,21,23,34$; its
bidirected edges are $01,02,03,13,14,24$. The only four-cycle is
$0\!-\!1\!-\!4\!-\!2\!-\!0$, corresponding to $O$. For each repeated colour,
exactly two bicoloured corners occur. They fix the listed adjacent variety
up to reversal, and every reverse requires a missing arc. These six
varieties are therefore exactly the covered targets. The private
tricoloured positions are $30,31,32,40$ for $X_0$, none for $X_1,X_2,X_3$,
$10,14,20,24$ for $O$, and $12,13$ for $Y$; every adjacent target also has
two private bicoloured positions. Each target needs one more copy than is
available. The collection fails and its labelled orbit has size 240.

\begin{theorem}[Fourteen-cube obstructions]\label{thm:classification14}
Every fourteen-cube obstruction is equivalent to a member of exactly one
of the following forms:
\begin{enumerate}
 \item the failing collections of Corollary~\ref{cor:extremal-family};
 \item the all-adjacent derangement family above;
 \item the mixed exception above.
\end{enumerate}
There are respectively $51,4,1$ symmetry classes, hence 56 in total.
Their distribution by number of varieties is
\[
\begin{array}{c|rrrr|r}
 \text{varieties}&3&4&5&6&\text{total}\\\hline
 \text{classes}&5&20&26&5&56\\
 \text{labelled collections}&840&4800&5280&510&11430.
\end{array}
\]
\end{theorem}
\begin{proof}
The preceding arguments prove failure of all three forms. In the first
form, there are 91 distributions satisfying the disjoint-repeat condition.
The only residual nonidentity symmetry sends
\[
 (x_2,x_3,x_4;y_1,y_2,y_4)
 \longmapsto (y_4,y_1,y_2;x_3,x_4,x_2).
\]
It fixes 11 distributions: $y_4=x_2$, $y_1=x_3$, $y_2=x_4$ and
$x_2x_4=0$. Burnside's lemma gives $(91+11)/2=51$ classes. The three
forms cannot be equivalent, since their opposite-repeat totals are
respectively four, zero and three.

Necessity is the finite part of the proof. The anchored Hall cover applies
to fourteen sources. Theorem~\ref{thm:multiplicity-five}, independently of
the fifteen-cube threshold, removes the eight roots with anchor
multiplicity six or seven. The remaining 349 roots use total fourteen.
Their interval certificates retain 81 explicit terminal count vectors.
In addition to the exclusion, refinement and splitting rules, the checker
accepts an exact singleton box only when its vector occurs in the terminal
list. Thus every admissible vector in every root belongs to that list.
Each terminal is carried to one of the 56 representatives by an explicit
common colour/reflection symmetry. Separate finite checks verify the
Hall deficiencies, the full symmetry group and pairwise inequivalence.
The retained data give the displayed orbit and support counts.
\end{proof}

In particular, any fourteen-cube collection using at least seven varieties solves. The all-adjacent obstructions
are exactly the derangement family.

The classification also explains the corner threshold by a short deletion
argument. Suppose an unsolvable fifteen-cube pool existed. Every deletion
would be a classified fourteen-cube obstruction. If all sources were
adjacent-repeat, every deletion would have profile $(3,3,3,3,1,1)$.
A variety occurring three times in one deletion had multiplicity three or four in the
original pool. Deleting from that variety in the first case, or from
another variety in the second, leaves multiplicity two or four, both
impossible.

Otherwise there is exactly one opposite-repeat variety $A$: two distinct
ones could both be retained by a deletion. There are adjacent sources as
well, since every classified obstruction has at most four opposite copies.
Deleting an adjacent copy shows that the multiplicity of $A$ is three or
four; deleting an $A$ shows that one less is also three or four. Hence it
is four. Deleting an $A$ gives the mixed exception, so the adjacent
multiplicities in the original pool are $(5,3,1,1,1)$. Delete from the
five-copy variety. The resulting obstruction has four opposite copies
and six varieties, whereas the first form has at most five varieties.
This final contradiction proves the fifteen-cube upper bound.


\begin{theorem}[Three-variety obstructions]\label{thm:three-variety-classification}
Up to this equivalence, the fourteen-cube obstructions using exactly three
varieties are precisely five collections $4A+5B_r+5C_s$. The varieties are those of Corollary~\ref{cor:extremal-family};
Class I is the original illustrated obstruction. Representatives are
\[
\begin{array}{c|ccccc}
 \text{class}&\mathrm{I}&\mathrm{II}&\mathrm{III}&\mathrm{IV}&\mathrm{V}\\\hline
 r&2&2&3&3&4\\
 s&1&4&1&2&2
\end{array}
\]
\end{theorem}
\begin{proof}
The other two forms in Theorem~\ref{thm:classification14} have six varieties.
A three-variety member of the first form must concentrate each five-copy
group on one variety, giving $4A+5B_r+5C_s$ with $r\ne s$.
Corollary~\ref{cor:extremal-family} proves failure for all seven such pairs.
A symmetry between normalized collections fixes the unique
four-copy variety $A$, hence fixes $0$ and preserves its two unordered
pairs. Without arc reversal it fixes both cuts, hence every vertex.
With reversal it exchanges the cuts and is $\sigma=(1\ 3)(2\ 4)$.
It sends $(r,s)$ to $(\sigma(s),\sigma(r))$, pairing $(2,1)$ with $(3,4)$
and $(3,2)$ with $(4,1)$, while fixing $(2,4),(3,1),(4,2)$.
This gives exactly the five inequivalent representatives displayed.
\end{proof}
This displayed proof uses the complete classification. The formal endpoint
instead has an independent proof without interval certificates
(Appendix~\ref{app:extremal}); the repository also retains the elementary
necessity argument in \path{paper/three-variety-necessity.tex}.
